\documentclass[a4paper, 12pt]{article}
\usepackage[portuges]{babel}
\usepackage[utf8]{inputenc}
\usepackage{amsmath}
\usepackage{indentfirst}
\usepackage{graphicx}
\usepackage{multicol,lipsum}

\usepackage{xcolor}
\usepackage{hyperref}
\hypersetup{
    colorlinks=true,
    % linkcolor=cyan,
    filecolor=magenta,      
    urlcolor=blue,
    citecolor=green,
}
\usepackage{titlesec}


\titleformat{\subsection}
  {\normalfont\normalsize\bfseries}
  {\thesubsection}{1em}{}
  [\vspace{-0.3em}\rule{\textwidth}{0.3pt}]

\begin{document}
%\maketitle

\begin{titlepage}
	\begin{center}
	
	%\begin{figure}[!ht]
	%\centering
	%\includegraphics[width=2cm]{c:/ufba.jpg}
	%\end{figure}

		\Huge{UDESC}\\
		\large{Departamento de Ciência da Computação}\\ 
		\vspace{15pt}
        \vspace{95pt}
        \textbf{\LARGE{Teoria dos Grafos: Tarefa 1}}\\

		\textbf{Representação Computacional de um grafo não direcionado através de uma Lista de Adjacências}
		%\title{{\large{Título}}}
		\vspace{3,5cm}
	\end{center}
	
	\begin{flushleft}
		\begin{tabbing}
			Alunos: Heitor Linhares Caetano e Deivid Veras Lourenço \\
			Professor: Gilmário Barbosa dos Santos \\
	      \setlength{\labelsep}{\leftmargin}
	\end{tabbing}
 \end{flushleft}
	\vspace{1cm}
	
	\begin{center}
		\vspace{\fill}
			 Maio\\
			 2026
			\end{center}
\end{titlepage}
%%%%%%%%%%%%%%%%%%%%%%%%%%%%%%%%%%%%%%%%%%%%%%%%%%%%%%%%%%%

% % % % % % % % %FOLHA DE ROSTO % % % % % % % % % %

\newpage
% % % % % % % % % % % % % % % % % % % % % % % % % %
\newpage
\tableofcontents
\thispagestyle{empty}

\newpage
\pagenumbering{arabic}
% % % % % % % % % % % % % % % % % % % % % % % % % % %

\section{Objetivo do Trabalho}

A atividade consistiu em implementar um grafo através da representação computacional
\textit{Lista de Adjacências} e que é carregado a partir de uma base de dados no
formato CSV.

A Lista de Adjacências é uma estrutura de dados composta por um vetor estático
de Nós (unidade básica de uma lista encadeada). No nosso caso, cada índice do
vetor representa um vértice diferente (v1, v2, ...), e cada lista associada a um
índice representa as adjacências do vértice correspondente.

\begin{figure}[h]
\centering
\begin{tabular}{c|l}
Vértice & Adjacentes \\ \hline
1 & 2, 3, 5 \\
2 & 1, 4 \\
3 & 1 \\
4 & 2, 5 \\
5 & 1, 4 \\
\end{tabular}
\caption{Representação de uma lista de adjacências. O vértice 1 é
conectado a 2, 3 e 5.}
\end{figure}

Foram implementadas funções que fornecem o valor máximo e mínimo de um vértice,
funções que determinam o tipo do grafo (simples ou multigrafo), além da quantidade
de laços e/ou arestas múltiplas que ocorrem nele e, como função principal, um
procedimento que determina quantos componentes conexos existem no grafo, bem como
sua distribuição e seus respectivos tamanhos.

\newpage
\section{Solução fornecida e sobre o projeto}

O projeto foi administrado em um \href{https://github.com/atomicpenguinz/TEG_1}{repositório}
do GitHub, e separado em arquivos diferentes para melhor organização do trabalho.
Um arquivo para as funções referentes à Lista de Adjacências, outro referente
a leitura e interpretação do arquivo CSV, outra para a interface de usuário e ainda
um arquivo header, que possui as definições e protótipos usados nos outros.

A seguir vai uma explicação simples sobre cada exigência:

\subsection{Leitura do CSV}
Primeiramente, é importante ressaltar que o arquivo enviado não é exatamente um
CSV, pois os valores ali estão separados por espaços, e não por vírgulas. Isso
não é um grande problema, visto que alterar o código para suportar esse "SSV"
seria trivial.

Mesmo assim, decidimos alterar o arquivo da base de dados a fim de transformá-lo
em um CSV. Se tratou de uma tarefa simples, necessitando apenas os seguintes comandos
em \textit{Bash}:

\begin{figure}[h]
\centering
	\begin{verbatim}
	$ tr ' ' ',' < teste2.csv > tmp.csv
	$ mv -f tmp.csv teste2.csv
	\end{verbatim}
	\vspace{-1.0cm}
	\caption{O comando \textit{tr} está trocando todas as aparições do caractere
	espaço por uma vírgula.}
\end{figure}

Antes de começarmos a ler os valores, precisamos determinar o tamanho do vetor
estático que abrigará as adjacências. Para isso, foi criada uma função que corre
por todos os números e acha o maior dentre eles. Esse número (+ 1) é então utilizado como
tamanho do vetor, representando o maior índice existente no grafo.

Um possível problema dessa implementação é a eventualidade de não existir algum
vértice entre 1 e o maior número, o que poderia significar que ele simplesmente
não existe, mas aqui é interpretado como um vértice de grau nulo. Em uma conversa
com o professor, o mesmo declarou nossa solução como válida.

Ao ler o arquivo, como estamos tratando de grafos não direcionados, para cada par
de valores \textit{a,b} precisamos adicionar um novo \textit{nó} às listas de
adjacências de ambos os vértices. No caso em que \textit{a = b}, essa inserção é
realizada apenas uma vez, a fim de evitar a duplicação de laços.

\subsection{Grau mínimo e máximo}

Aqui a implementação foi relativamente direta. Nos dois casos, todas as listas
têm seus tamanhos contados através de um laço de repetição e o menor/maior número
é retornado no final.

\subsection{Quantidade de laços e arestas múltiplas}

A existência de laços e/ou arestas múltiplas em um grafo confirma que este é um
multigrafo. Existem várias abordagens possíveis para a contagem de valores repetidos,
e a complexidade de cada uma foi levada em conta para a nossa escolha final.

Inicialmente, a ideia foi de fazer a contagem através de um vetor de frequência.
Nessa estratégia, para cada vértice adjacente na lista, incrementa-se a posição
correspondente em um vetor auxiliar. Ao final da leitura, valores maiores que 1
significam a presença de múltiplas arestas. Isso nos permite contabilizar duplicatas
e laços de maneira linear.

Essa abordagem tem, para cada vértice, complexidade de tempo \textit{O(g + V)},
onde g é o grau do vértice e V o número total de vértices do grafo. Isso porque
um vetor de tamanho V era inicializado e percorrido em cada vértice, o que significa
que a complexidade de espaço também era \textit{O(V)}.

Em grafos esparsos (como
o analisado), isso significa que ocorre um desperdício de memória, pois grande
parte do vetor nem chega a ser usada.
Como o procedimento é executado para cada vértice, a complexidade total da função
era \(O(V^2 + E)\), onde \(E\) é o número de arestas.

\subsection{Determinação de componentes conexos}
\newpage

\section{Sobre o ambiente de desenvolvimento e compilação}
Como explicado anteriormente, um repositório de código-fonte foi instaurado no
GitHub para uma melhor colaboração.

Como sistema operacional, foram utilizados tanto o Windows (Deivid) quanto GNU/Linux
(Heitor). Mesmo com ambientes diferentes, miramos em ter uma \textit{codebase}
compatível, sem usar funções não portáteis, como \textit{system()}.

Sobre os ambientes de desenvolvimento, Deivid usou o \href{https://code.visualstudio.com/}
{Visual Studio Code}, enquanto Heitor usou o editor de texto \href{https://helix-editor.com}
{Helix}. Como este não usou uma IDE completa, algumas ferramentas de linha de comando
foram usadas para emular funcionalidades comuns em outros ambientes, especificamente
\href{https://invisible-island.net/cproto/cproto.html}{cproto} para a geração automática
de protótipos e \href{https://astyle.sourceforge.net/}{astyle} como formatador de código.

Como as demais ferramentas, também foi utilizada a linha de comando para compilar
o projeto, em que foi adotado o compilador de C do \textit{GNU (GCC)}.

O comando para compilar é:

\verb|$ gcc ls_adj.c csv.c main.c -o grafo|

e para iniciar o programa, simplesmente execute:

\verb|$ ./grafo|

\newpage
\section{Análise dos Resultados}
\newpage

\section{Conclusões}
\end{document}
